Ze bleef in de vulva
want baby Maria

wantrouwde de wereld.
Ze was warm gezeteld

in moeders baarmoeder.
De veiligste hoeder

was vriendelijk op zich.
Buiten was vijandig.

Al in haar eerste jaar
vergrootte alles haar

basis van wantrouwen.
De velerlei vrouwen,

die kwamen als hoeders
waaronder zoogmoeders

en toen kindermeiden,
wist hechting te mijden.

De verzorging was kil.
Niet slechts voor haar bestwil,

maar voor een inkomen
waren ze gekomen.

Men bleek onbetrouwbaar
en vaak onuitstaanbaar.

Met Erik Erikson
als informatiebron

zou men “ik-ben-waardig”
of “de-mens-is-aardig”

nergens bij Maria
in levensstadia

terug kunnen vinden.
Ze kon zich niet binden.

Elk gebaar werd gebumpt.
Ze werd direct gedumpt

in het kloosterleven
toen Maria zeven

jaar oud was geworden.
Haar grootouders morden,

maar haar vader hield stand.
Hij deed haar van de hand.

Haar synchroniciteit
raakte Maria kwijt.

Weer het bekende lied.
Zorgen maken hielp niet.

Het waren gedachten,
die kwamen als klachten.

Ze wilden afleiden
van het ware lijden

wat in het hier en nu
echt stond op het menu.
Maria moest eten
voor een dieper weten.

Maria was waakzaam,
maar dat bleek niet heilzaam.

Als ze niet alert was
dan faalde haar harnas

gesmeed uit perfectie,
voor Maria’s protectie.

‘The fear of missing out’.
De paniek bij een fout.

Al zag ze niet alles
en gaf dat hommeles,

bleef ze acht volhouden.
In de gaten houden

alsof je nooit iets mist.
Had ze zich toch vergist

dan groeide wantrouwen
alsook het vertrouwen

in ‘op de hoede zijn’.
“Het leven brengt slechts pijn

als je niet goed oplet”,
was Maria’s verzet

tegen de overgave.
Spanning haar opgave.

Op een donkere schaal
midden in de eetzaal

lag een kleverig blok
vlak naast een kaneelstok.

Het glinsterde haar toe.
Ze wantrouwde haar moe

dat na de visite
van haar favoriete

lekkers - na marsepein -
nog wat over zou zijn.

De turrón of noga
ging deels in Maria,

voor de gasten kwamen.
Ze kon zich wel schamen,

maar anders kreeg ze niets.
“Eet dan tenminste iets”

dacht ook Pipi Langkous.
De zoete honingsaus

over amandelen
vergaf haar handelen.

Ook God leek dat te doen.
Haar moeder gaf haar toen

een eigen snoepschaaltje
met het fout verhaaltje

dat ze braaf was geweest.
Het maal rook als een feest.

Vooral de marsepein
vond Maria heel fijn.

Tot ze even later
spijt kreeg van haar flater
omdat ze moest kotsen.
Want hard en zacht botsen

net als vet en honing,
plakken versus smering.

Moeder dacht even na.
“Nooit snoep voor Maria!”,

was toen haar conclusie
en de koks instructie.

Menigeen grinnikte
toen Maria predikte:

“Bangheid uit wantrouwen
leidt vaak tot vertrouwen

als het mag uitbreken
en mag openbreken.

De muren vallen weg
voor een verbindingsweg

naar andere mensen.
Zonder de pijngrenzen

wordt je niet gehinderd.
Liefde niet geminderd

door schijnonkwetsbaarheid,
creëert verbondenheid!”

“Spanning wordt frustratie”,
aldus hoorde Emmy,

“als we machteloos zijn.
Onzekerheid doet pijn.”

Een eenling vermoedde
wat in haar hoofd broedde:

“Ik die de bijbel goed ken,
geloof : ‘als ik bang ben,

dan vertrouw ik op U!’
Waar is dat geloof nu?”

Meest ruilde wantrouwen
aldoor met vertrouwen

afhankelijk of vrees
zijn bedekking verrees.

“Zelfs die van je houden,
kan je niet vertrouwen,

want we zijn allen bang!”,
predikte ze al lang.

“En als we dat zagen,
dan kon ik het wagen

mensen te vertrouwen”
Maria kon rouwen,

want die notie was erg.
Nergens was een herberg

waar je safe kon schuilen.
Waar je vrij kon huilen

en altijd werd getroost.
Altijd werd geliefkoosd

en dus kon genieten
wat ouders nalieten.

Als ze vertrouwen greep
en haar hand samen kneep

dan glipte de hoop weg.
Het knijpen gaf uitleg

aan haar radeloosheid.
Haar geloof was al kwijt.

In haar bangigheid gaan,
heeft ze nooit echt gedaan.

Vertrouwen werd een doel,
in plaats van een gevoel,

dat volgt uit doorvoelen.
Angst laten afkoelen.

Nooit straalde Maria,
licht als een gloria,

verbonden met het nu.
Slechts als individu

ondervond ze haar kern.
De wereld was extern

verpakt in verhalen,
die zich steeds herhalen.

Het hart werd niet geraakt.
De liefde werd gestaakt.

Geen inzicht werd ontwaakt.
Niets werd helder gemaakt.

Reed Mullah Nasruddin
niet ooit een dorpje in

zittend op een ezel?
Zat de wijze wezel

niet achterstevoren
kijkend naar diens oren?

Het dorp vroeg uit angst niets,
ook al vond men zoiets

natuurlijk ongewoon.
Hij bleef een wijs persoon.

Vroeg toen niet een klein kind
“Wijze Hodja, ik vind

uw positie raar!
Uw ogen kijken naar

wat al achter u ligt?”
Mullah keek naar het wicht

en zei toen bovendien
“Dat heb je goed gezien!

Dat is wat zeker is.
Wat ik zelf niet beslis,

is hoe de toekomst loopt.
Wat deze dienaar hoopt

is dat de ezel weet
welke steen, grond of spleet

beiden laat struikelen.
Ik zal niet duikelen

als ik hem blind vertrouw
en niet trek aan het touw.”

Zoals Kierkegaard zei:
“Slechts leven al voorbij

kan begrepen worden.”
Het zijn aanwijsborden

naar de vage uitkomst
van ieder zijn toekomst.

Kon God argwaan leren,
door te camoufleren

wat hij werkelijk was?
Emmy als alias?

Maria als eerste,
de origineelste?

In hun imperfectie
vond hij zijn perfectie.

Zonder die beperking
was er geen ontdekking.

“Heb ik als mens een doel?
Is menszijn als gevoel

soms God die zich ervaart?
Is dat mijn leven waard?

Elke verandering
is dan een verrassing!

Elk wantrouwen komt dan
van twijfel bij Gods plan.

Dus van onwetendheid
van Gods voorzienigheid.

Want Jezus zei mede
tijdens zijn bergrede

dat God je leven borgt
en goed voor alles zorgt,

als ik Lucas 12 snap”,
vertolkte Emmy knap.

Gold dat ook voor Emmy
in deze rijmfictie?

Kon Emmy zeker zijn
over haar vaarwelzijn?
